\section{Research Gap and Positioning of Motivo}\label{sec:motivo-position}

The domain-specific languages surveyed above split into two broad families.
Live-coding systems prioritize interactive performance and pattern transformation, streaming events to a synthesis
backend rather than producing a self-contained offline artifact~\cite{Aaron_2016,McLean_Wiggins_2010,Roos_McLean_2023,
Du_Bois_Ribeiro_2019}.
Real-time audio languages offer precise control over synthesis and timing, but at an abstraction level oriented
toward signal graphs and concurrent processes rather than hierarchical song
structure~\cite{Wang_Cook_2003,McCartney_2002}.
None of the surveyed DSLs treats offline, deterministic compilation to a portable symbolic file as its primary design
goal, nor do they combine typed imperative source with a zero-runtime expansion model.

From this analysis, three gaps among existing musical DSLs emerge.
First, most target live performance or synthesis, leaving offline deterministic compilation underexplored as a primary
design goal.
Second, few systems bridge low-level audio programming and high-level compositional structure with a zero-runtime
compilation model.
Third, composition and execution are often conflated, so musical intent remains tied to performance-time scheduling.

Motivo addresses these gaps by providing a typed, imperative source language for offline symbolic composition;
a compiler that validates structure and expands it deterministically into a flat event timeline; and a Standard MIDI
File as the portable output artifact.
Patterns and loops compress repetitive structure in the source program while the lowering stage materializes the full
event stream, as introduced in the compression argument above.

Having located Motivo against these systems, the next chapter turns from \emph{why} such a language is warranted to
\emph{what} it actually is: Chapter~\ref{ch:language-design} specifies Motivo's conceptual model, grammar, type system,
and temporal semantics, making concrete the typed, offline, zero-runtime design just argued for.
